\documentclass[12pt]{article}
\PassOptionsToPackage{table}{xcolor}
\usepackage[utf8]{inputenc}
\usepackage[T1]{fontenc}
\usepackage[italian]{babel}
\usepackage{multicol}
\usepackage{enumitem}
\usepackage{array}
\usepackage{longtable}
\usepackage{booktabs}
\usepackage{xcolor}
\usepackage{graphicx}
\usepackage{tikz}
\usetikzlibrary{
    shapes.geometric,
    arrows.meta,
    positioning,
    calc,
    patterns,
    decorations.pathmorphing
}
\usepackage[margin=2cm]{geometry}
\usepackage{amsmath,amssymb}
\usepackage{siunitx}
\setlength{\parindent}{0pt}
\definecolor{headercolor}{RGB}{220,220,220}
\newcounter{quesito}
\setcounter{quesito}{0}
\newcommand{\nuovoquesito}{\stepcounter{quesito}\textbf{Domanda \arabic{quesito}.}}
\newcommand{\itemdomanda}[2]{%
    \par\noindent
    \begin{minipage}[t]{\linewidth}
        \nuovoquesito \\ #1
        \begin{enumerate}[label=(\Alph*), nosep, leftmargin=*]
            #2
        \end{enumerate}
    \end{minipage}
    \vspace{10pt}
}
\begin{document}
\section*{Simulazione}
\hfill Data: 17 apr 2026

\itemdomanda{Un mattone pesa un chilo più mezzo mattone. Quanto pesa un mattone?}{
    \item \SI{1.5}{kg}
    \item \SI{2}{kg}
    \item \SI{2.5}{kg}
    \item \SI{3}{kg}
    \item \SI{1}{kg}
}

\itemdomanda{Il figlio di mio padre è il padre di mio figlio. Chi sono io?}{
    \item Mio nonno
    \item Mio zio
    \item Mio padre
    \item Io stesso
    \item Mio figlio
}

\itemdomanda{Completare la serie: $1, 8, 27, 64, \dots$}{
    \item 100
    \item 81
    \item 125
    \item 144
    \item 216
}

\itemdomanda{Completare la serie: $2, 6, 12, 20, 30, \dots$}{
    \item 36
    \item 40
    \item 42
    \item 50
    \item 44
}

\newpage
\section*{Appendice: risposte corrette}
\begin{longtable}{@{}r>{\raggedright\arraybackslash}p{0.14\linewidth}>{\raggedright\arraybackslash}p{0.20\linewidth}>{\raggedright\arraybackslash}p{0.42\linewidth}c@{}}
\toprule
\textbf{N.} & \textbf{Materia} & \textbf{Argomento} & \textbf{Titolo} & \textbf{Risp.} \\
\midrule
\endfirsthead
\toprule
\textbf{N.} & \textbf{Materia} & \textbf{Argomento} & \textbf{Titolo} & \textbf{Risp.} \\
\midrule
\endhead
\midrule
\multicolumn{5}{r}{\footnotesize\textit{Continua}} \\
\endfoot
\bottomrule
\endlastfoot
1 & Logica & Logica Analitica & Ordinamento pesi & B \\
2 & Logica & Logica Analitica & Relazioni familiari & D \\
3 & Logica & Logica Numerica & Successione cubica & C \\
4 & Logica & Logica Numerica & Successione quadratica & C
\end{longtable}

\end{document}
